% Введение: здесь обычно формулируются актуальность, цель и практическая значимость.
% Студент может адаптировать формулировки под свою предметную область.
Современная образовательная среда развивается в логике цифровой трансформации, при которой учебный процесс перестает быть только последовательностью очных занятий и становится непрерывной системой взаимодействия между студентом, преподавателем и цифровой платформой. В университетских и корпоративных курсах накапливаются детальные данные о поведении обучающихся: попытки решения задач, длительность прохождения модулей, частота повторного обращения к материалам, характер типичных ошибок и динамика восстановления после неудачных шагов. Если раньше эти данные использовались преимущественно для формальной отчетности, то сегодня они могут служить основой для принятия педагогических решений в реальном времени \autocite{dumont2023personalized,zawacki2019}. Именно поэтому адаптивное обучение рассматривается не как локальная технологическая опция, а как стратегическое направление развития цифрового образования.

Рост объема образовательных данных сам по себе не гарантирует повышения качества обучения. В традиционной модели преподаватель вынужден агрегировать большой поток сигналов вручную, интерпретируя их на уровне экспертных суждений и опыта. Такой подход остается ценным, но в масштабных курсах он ограничен ресурсами времени и внимания. Практика показывает, что при фиксированном учебном плане часть группы систематически не успевает за темпом, а часть теряет мотивацию из-за недостатка интеллектуальной сложности. Возникает противоречие между массовым форматом реализации дисциплины и потребностью в индивидуальных траекториях, где темп, уровень сложности и тип обратной связи подбираются под конкретного обучающегося \autocite{kornilov2022adaptive,osmolovskaya2018}.

Теоретические основания персонализации образования сформировались задолго до появления современных платформ и алгоритмов. В работах Дж. Дьюи и последующей педагогической традиции акцент делался на деятельностном характере обучения и необходимости учитывать интересы, подготовку и контекст обучающегося \autocite{dewey2000,pedencyclopedia1972}. На практическом уровне эти идеи воплощались через дифференцированные задания, вариативные маршруты и индивидуальные консультации. Однако в условиях массового высшего образования ручная реализация такой модели затруднена: преподавателю требуется одновременно отслеживать множество параметров по всей академической группе. Цифровые инструменты и методы машинного обучения открывают возможность частично автоматизировать этот контур, сохраняя роль преподавателя как носителя педагогической ответственности.

Существенным драйвером развития цифровых образовательных платформ стал период массового перехода к дистанционным форматам, когда вопрос качества взаимодействия и персональной поддержки студентов стал критически важным. После этого этапа многие практики, первоначально воспринимавшиеся как вынужденные, закрепились в регулярном образовательном процессе. В частности, использование трекеров активности, модульных заданий и автоматизированной диагностики перестало быть исключением. Вместе с этим усилился запрос на прозрачные механизмы персонализации, где каждый следующий шаг курса объясняется не только статистической моделью, но и интерпретируемым педагогическим основанием. В инженерном представлении это означает, что сущности вроде \texttt{learning\_event}, \texttt{skill\_id} и \texttt{mastery\_score} должны быть не просто техническими полями в базе, а элементами единой образовательной логики.

На уровне прикладных решений рынок демонстрирует широкий спектр подходов: от платформ с минимальной адаптацией темпа до систем, использующих прогнозирование успеваемости, оценку риска отчисления и рекомендацию контента. Международные и российские EdTech-продукты подтверждают, что персонализация повышает вовлеченность и удержание, особенно в длинных программах, где критичны своевременная обратная связь и корректный подбор сложности \autocite{bicknell2023,mon2023rl}. Одновременно исследования отмечают риски: непрозрачность модели, смещение рекомендаций, зависимость качества результатов от полноты входных данных и качества разметки педагогических сценариев. Это создает требование к архитектуре системы, в которой алгоритмический слой проектируется вместе с механизмами объяснения и экспертного контроля \autocite{said2024,sok2024llm}.

В рамках данной работы адаптивное обучение рассматривается как управляемый процесс, где на каждом шаге необходимо сопоставить текущее состояние студента и множество доступных педагогических действий. В терминах системы это соответствует циклу: сбор данных, обновление модели знания, выбор следующего задания, фиксация результата и переоценка стратегии. Практически цикл может быть представлен как последовательность операций над объектами \texttt{student\_profile}, \texttt{task\_bank}, \texttt{policy\_state} и \texttt{feedback\_trace}. Такой подход позволяет связать педагогические категории с инженерными интерфейсами, не теряя содержательной части процесса.

Актуальность исследования определяется сочетанием трех факторов. Во-первых, цифровая инфраструктура образования уже предоставляет достаточный объем данных для построения персонализированных траекторий. Во-вторых, существуют апробированные алгоритмические подходы к оценке знания и последовательному принятию решений в условиях неопределенности. В-третьих, учебная практика нуждается в системах, которые дают не только "рекомендацию", но и внятное объяснение причин выбора конкретного шага. Для этого требуется интеграция статистических моделей, методов машинного обучения и инструментов генерации интерпретаций в едином программном контуре.

Целью исследования является разработка и демонстрация воспроизводимого подхода к адаптивному обучению, в котором алгоритмы машинного обучения используются для персонализации образовательного контента и темпа прохождения материала. Достижение цели предполагает не только техническую реализацию прототипа, но и формализацию критериев качества персонализации, чтобы результаты могли быть оценены с позиций педагогической практики.

Для достижения поставленной цели в работе решаются следующие взаимосвязанные задачи. Первая задача состоит в анализе существующих концепций адаптивного обучения и определении места методов машинного обучения в современной образовательной среде. Вторая задача связана с формализацией архитектуры интеллектуального компонента системы, включая модель представления знания и механизм выбора следующего шага обучения. Третья задача предполагает проектирование программного прототипа, в котором выделены роли студента, методиста и эксперта, а контур рекомендаций поддерживает как автоматические, так и экспертно корректируемые решения. Четвертая задача направлена на демонстрацию функциональных возможностей прототипа и интерпретацию результатов на основе метрик качества рекомендаций.

Научно-практическая значимость работы заключается в том, что предложенный шаблон реализации позволяет переносить разработанный подход в новые дисциплины без полного пересмотра архитектуры. За счет явной конфигурации параметров и модульного деления компонентов возможно адаптировать систему под разные учебные сценарии, изменяя набор навыков, правила диагностики и стратегию выбора заданий. Для эксплуатации это означает, что конфигурация уровня \texttt{epsilon\_start}, \texttt{reward\_weights} и \texttt{mastery\_threshold} может регулироваться предметным экспертом без вмешательства в ядро платформы.

Структурно работа включает введение, две главы, заключение и список источников. В первой главе раскрывается теоретический и методический контур персонализации обучения, во второй главе описывается программная реализация, пользовательские сценарии и результаты тестирования. В заключении формулируются основные выводы и направления дальнейшего развития, включая расширение корпуса данных, интеграцию с внешними LMS и проведение длительных педагогических экспериментов на реальных образовательных потоках.
